\appendix
\chapter{Lean Setup and Conventions}
\label{app:lean-setup}

\section{Minimal Project Setup}

The setup here is identical in shape to \textit{Proof and
Countermodel}'s appendix, since both books target the same Lean~4
toolchain and the same one \texttt{mathlib} dependency; only the
project name differs.

\texttt{lean-toolchain}:
\begin{lean}
leanprover/lean4:v4.11.0
\end{lean}
\texttt{lakefile.toml}:
\begin{lean}
name = "inference_under_constraint"
defaultTargets = ["InferenceUnderConstraint"]

[[require]]
name = "mathlib"
scope = "leanprover-community"

[[lean_lib]]
name = "InferenceUnderConstraint"
\end{lean}
As before, \texttt{lake update} followed by \texttt{lake build}
fetches and compiles \texttt{mathlib} once; every listing in this
book can then be pasted into a \texttt{.lean} file under
\texttt{InferenceUnderConstraint/} and checked directly.

\section{Why \texttt{K3} Is an Inductive Type, Not \texttt{Bool}}

\textit{Proof and Countermodel} builds almost everything on
\texttt{Bool}, since a two-valued register is exactly what that
book's classical material needs. This book cannot do the same: its
first chapter needs a genuine \emph{third} value, and its second
needs a small, named set of possible worlds. Both are the same kind
of move, and it is worth naming the move once here rather than
re-justifying it in each chapter.

\begin{lean}
inductive K3 where
  | tt | ff | unknown
  deriving DecidableEq, Repr
\end{lean}
An \texttt{inductive} type with explicitly named constructors gives
\texttt{K3.tt}, \texttt{K3.ff}, and \texttt{K3.unknown} as three
distinct, self-describing values, and the \texttt{deriving
DecidableEq} clause is what lets \texttt{decide} enumerate over them
exactly as it enumerates over \texttt{Bool} in the first book — the
finiteness that makes \texttt{by decide} work is a property of
\texttt{DecidableEq} plus a finite constructor list, not a property
special to \texttt{Bool} itself.

The alternative — encoding the three values as, say, \texttt{Fin 3}
and matching on numeric literals (\texttt{0}, \texttt{1}, \texttt{2})
— was deliberately avoided. Literal patterns over \texttt{Fin n} are
sensitive to how the underlying numeral is elaborated and have been
known to fail to match exhaustively across Lean/mathlib versions in
ways that are silent until the version changes; named constructors
sidestep this entirely, since \texttt{match} on an inductive type's
own constructors is checked for exhaustiveness by the type itself,
not by numeral coincidence. Chapter~2's Kripke frame
(\texttt{World} with constructors \texttt{w0}, \texttt{w1},
\texttt{w2}) uses the identical pattern for the same reason.

\begin{thebibliography}{9}

\bibitem{priest2008}
G.~Priest,
\emph{An Introduction to Non-Classical Logic: From If to Is}, 2nd ed.,
Cambridge University Press, 2008.

\bibitem{priest2006}
G.~Priest,
\emph{In Contradiction: A Study of the Transconsistent}, 2nd ed.,
Oxford University Press, 2006.

\bibitem{blackburn2001}
P.~Blackburn, M.~de~Rijke, and Y.~Venema,
\emph{Modal Logic},
Cambridge Tracts in Theoretical Computer Science \textbf{53},
Cambridge University Press, 2001.

\bibitem{kleene1952}
S.~C. Kleene,
\emph{Introduction to Metamathematics},
Van Nostrand, 1952.

\end{thebibliography}
